\section{\uppercase{Conclusions}}
\label{sec:conclusion}

We asked what limits species-level reading of sticky traps and which lever moves that limit. Along a ladder of five cumulative training sets, species accuracy rises from 56.1\% to 82.2\%, with the largest gains on objects below 32\,px and insect-only plates and part of the look-alike confusion shrinking; but random subsamples of the final training set of the size of L1 and L2 reach the same accuracy, so from L1 to L2 the rise is mostly the number of images, not targeting, but not the L0 deficit (a random 103-image subsample is 19\,pp better). Identification did not move with the classifier: with our recipe, a detect-then-classify cascade does not improve on the single-stage detector in any of 15 configurations (nor with species-only relabelling, late fusion or other crop margins), ground-truth-crop accuracy saturates near 83\% for five classifiers (four backbones), and the residual errors are mostly confident and concentrated in MD/MV/MC (only 4.9\% are wrong for all fifteen classifiers), so this is a plateau, not a proven information limit; label noise or the crop recipe cannot be excluded. Tiny-object methods from aerial imagery bring no gain detectable here: none of eight variants improves species accuracy over a same-campaign control with three seeds (all differences between $-3.8$ and $+0.3$\,pp except NWD with RFLA at $-8.8$\,pp), and none raises the recall of objects below 32\,px. For practitioners of trap-based monitoring, our results suggest that, among the levers tested, more annotated plates helped most (whether targeting adds beyond quantity remains open), and that investigating the remaining confusion within the housefly/blowfly/flesh-fly family as a possible labeling and imaging problem, which expert re-annotation could settle.

\paragraph{Limitations.} One facility and one test split (73 images; 7 crowded and 6 insect-only plates), whose 95\% intervals are about $\pm2$\,pp wide, so effects smaller than that cannot be resolved; three seeds; the cascade experiments use the final dataset only. Cross-trap-colour evaluation~\cite{ong2025} was not done; the split rule by plate is not documented, so leakage between images of the same plate cannot be excluded; RQ3 and part of the equal-size runs used a different GPU than RQ1/RQ2; checkpoints were selected on VAL, so no never-consulted hold-out exists. The tiny-object variants were tuned only for NWD (three settings, one with a single seed), so this is ``no gain detected here'', not proof of none; the NWD+RFLA drop stems from unstable seeds. Generalization beyond this facility is untested and the code and configurations will be released but the images will not, so the ladder cannot be reproduced independently.

\paragraph{Future work.} Out-of-distribution rejection for novel species~\cite{opensect2025}, expert re-annotation of the confusable pairs, a size control that separates label quality from quantity at the earliest lot, edge deployment.
